\documentclass[11pt]{article}
\usepackage[margin=0.85in]{geometry}
\usepackage{amsmath,amssymb}
\usepackage{xcolor,listings}
\usepackage{hyperref}
\hypersetup{colorlinks=true,linkcolor=blue}
\lstset{language=Python,basicstyle=\ttfamily\small,columns=fullflexible,
 keepspaces=true,showstringspaces=false,backgroundcolor=\color{black!3},
 frame=single,rulecolor=\color{black!15},aboveskip=7pt,belowskip=7pt}
\setlength{\parindent}{0pt}
\setlength{\parskip}{5pt}
\newcommand{\ket}[1]{\lvert #1\rangle}
\newcommand{\bra}[1]{\langle #1\rvert}
\newcommand{\mtwo}[4]{\begin{pmatrix}#1&#2\\#3&#4\end{pmatrix}}
\newcommand{\vfour}[4]{\begin{pmatrix}#1\\#2\\#3\\#4\end{pmatrix}}
\newcommand{\heading}[1]{\begin{center}\large\textit{#1}\end{center}}
\begin{document}

\heading{Matrix Product States: Construction and Storage}
\textit{Examples 05 and 06}

\subsection*{1. Three-spin coefficients}
Begin with
\begin{equation}
\ket{\psi}=\sum_{ijk}C_{ijk}\ket{ijk},\qquad
\sum_{ijk}|C_{ijk}|^2=1.
\end{equation}
For the GHZ state,
\begin{equation}
\ket{\psi}=\frac{\ket{000}+\ket{111}}{\sqrt2},\qquad
C_{000}=C_{111}=\frac1{\sqrt2},
\end{equation}
and all other coefficients vanish. Here $\ket0=\ket\uparrow$ and $\ket1=\ket\downarrow$.

\subsection*{2. First SVD: $1\,|\,23$}
Group the last two spins into one column index:
\begin{equation}
C_{ijk}\xrightarrow{\text{reshape}}M_{i,\,2j+k},\qquad
M=\begin{pmatrix}C_{000}&C_{001}&C_{010}&C_{011}\\
C_{100}&C_{101}&C_{110}&C_{111}\end{pmatrix}
=\frac1{\sqrt2}\begin{pmatrix}1&0&0&0\\0&0&0&1\end{pmatrix}.
\end{equation}
Factor the matrix:
\begin{align}
M&=U^{[1]}\Sigma^{[1]}V^{[1]\dagger},\\
M_{iq}&=\sum_\alpha U^{[1]}_{i\alpha}s^{[1]}_\alpha
(V^{[1]\dagger})_{\alpha q}.
\end{align}
Choose the first tensor and the remainder as
\begin{equation}
A^{[1]}_{i\alpha}=U^{[1]}_{i\alpha},\qquad
R_{\alpha q}=s^{[1]}_\alpha(V^{[1]\dagger})_{\alpha q},
\end{equation}
giving
\begin{equation}
C_{ijk}=M_{i,\,2j+k}=\sum_\alpha A^{[1]}_{i\alpha}R_{\alpha,\,2j+k}.
\end{equation}
The corresponding Schmidt form is
\begin{align}
\ket\psi&=\sum_\alpha s^{[1]}_\alpha\ket{a_\alpha}_1\ket{b_\alpha}_{23},\\
\ket{a_\alpha}_1&=\sum_i U^{[1]}_{i\alpha}\ket i,
\qquad\ket{b_\alpha}_{23}=\sum_{jk}(V^{[1]\dagger})_{\alpha,\,2j+k}\ket{jk}.
\end{align}
The index $\alpha$ labels Schmidt terms. For GHZ, there are two, both with singular value $1/\sqrt2$. One valid factorization is
\begin{equation}
A^{[1]}=I_2,\qquad R=\frac1{\sqrt2}\begin{pmatrix}1&0&0&0\\0&0&0&1\end{pmatrix}.
\end{equation}

\newpage
\subsection*{3. Second SVD: $12\,|\,3$}
Regroup the remainder so $(\alpha,j)$ labels rows and $k$ labels columns:
\begin{equation}
R_{\alpha,\,2j+k}\xrightarrow{\text{reshape}}M^{[2]}_{2\alpha+j,\,k},\qquad
M^{[2]}=\begin{pmatrix}R_{00}&R_{01}\\R_{02}&R_{03}\\R_{10}&R_{11}\\R_{12}&R_{13}\end{pmatrix}
=\frac1{\sqrt2}\begin{pmatrix}1&0\\0&0\\0&0\\0&1\end{pmatrix}.
\end{equation}
Apply SVD and separate the row index again:
\begin{align}
M^{[2]}_{2\alpha+j,\,k}
&=\sum_\beta U^{[2]}_{2\alpha+j,\,\beta}s^{[2]}_\beta
(V^{[2]\dagger})_{\beta k},\\
A^{[2]}_{\alpha j\beta}&=U^{[2]}_{2\alpha+j,\,\beta},
\qquad A^{[3]}_{\beta k}=s^{[2]}_\beta(V^{[2]\dagger})_{\beta k}.
\end{align}
Substituting into the first factorization gives
\begin{align}
R_{\alpha,\,2j+k}&=\sum_\beta A^{[2]}_{\alpha j\beta}A^{[3]}_{\beta k},\\
C_{ijk}&=\sum_{\alpha\beta}A^{[1]}_{i\alpha}A^{[2]}_{\alpha j\beta}A^{[3]}_{\beta k}.
\end{align}
The tensor shapes are
\begin{equation}
A^{[1]}:(2,\chi_1),\qquad A^{[2]}:(\chi_1,2,\chi_2),\qquad A^{[3]}:(\chi_2,2).
\end{equation}
Since $A^{[1]\dagger}A^{[1]}=I$, the retained left states are orthonormal. The second SVD therefore gives the Schmidt coefficients across $12\,|\,3$.

\subsection*{4. Explicit GHZ tensors}
One compatible choice is
\begin{equation}
A^{[1]}=I_2,\quad
A^{[2]}_{:,0,:}=\mtwo{1}{0}{0}{0},\quad
A^{[2]}_{:,1,:}=\mtwo{0}{0}{0}{1},\quad
A^{[3]}=\frac{I_2}{\sqrt2}.
\end{equation}
The colon includes all values of that index. Fixing the physical spins leaves ordinary matrix multiplication:
\begin{align}
C_{ijk}&=A^{[1]}_{i,:}A^{[2]}_{:,j,:}A^{[3]}_{:,k},\\
C_{000}&=\begin{pmatrix}1&0\end{pmatrix}\mtwo{1}{0}{0}{0}
\frac1{\sqrt2}\begin{pmatrix}1\\0\end{pmatrix}=\frac1{\sqrt2},\\
C_{001}&=\begin{pmatrix}1&0\end{pmatrix}\mtwo{1}{0}{0}{0}
\frac1{\sqrt2}\begin{pmatrix}0\\1\end{pmatrix}=0.
\end{align}
In general, $C_{ijk}=\delta_{ij}\delta_{jk}/\sqrt2$, recovering $(\ket{000}+\ket{111})/\sqrt2$. SVD factors can differ within degenerate Schmidt subspaces; their contraction gives the same state.

\newpage
\subsection*{5. Bond dimensions and storage}
Keeping all nonzero singular values gives exact bond dimensions
\begin{equation}
\chi_1=\operatorname{rank}C_{i,(jk)},\qquad
\chi_2=\operatorname{rank}C_{(ij),k}.
\end{equation}
For three spins, the entry counts are
\begin{align}
N_{\rm full}&=2^3=8,\\
N_{\rm MPS}&=(2\chi_1)+(\chi_1\,2\chi_2)+(\chi_2\,2).
\end{align}
For GHZ, $\chi_1=\chi_2=2$, giving $N_{\rm MPS}=16$. With only three spins, the MPS stores more entries than the full vector.

For $N$ spins, successive SVDs give
\begin{equation}
C_{i_1\cdots i_N}=\sum_{\alpha_1\cdots\alpha_{N-1}}
A^{[1]}_{i_1\alpha_1}A^{[2]}_{\alpha_1i_2\alpha_2}\cdots
A^{[N]}_{\alpha_{N-1}i_N}.
\end{equation}
With $\chi_0=\chi_N=1$, the stored entries are
\begin{equation}
N_{\rm MPS}=\sum_{n=1}^N2\chi_{n-1}\chi_n,\qquad N_{\rm full}=2^N.
\end{equation}
For GHZ, each internal bond has dimension two:
\begin{equation}
N_{\rm MPS}=4+8(N-2)+4=8N-8\qquad(N\geq2).
\end{equation}
The MPS count grows linearly with $N$, while the full-vector count doubles with each added spin. This advantage depends on the bond dimensions remaining small.

\subsection*{Why SVD}
At one split, retain the largest $\chi$ singular values:
\begin{equation}
M=\sum_\alpha s_\alpha u_\alpha v_\alpha^\dagger
\quad\longrightarrow\quad
M_\chi=\sum_{\alpha=1}^\chi s_\alpha u_\alpha v_\alpha^\dagger.
\end{equation}
The squared error in the matrix entries is
\begin{equation}
\|M-M_\chi\|_F^2
=\sum_{ij}|M_{ij}-(M_\chi)_{ij}|^2
=\sum_{\alpha>\chi}s_\alpha^2.
\end{equation}
The truncated SVD minimizes this error among rank-$\chi$ approximations. Removing zero singular values preserves the state; removing small ones gives an approximation with known error at that split. The code uses a tolerance to identify numerical zeros.

Repeated truncations do not necessarily give the globally best MPS at fixed bond dimensions. They apply this optimal matrix approximation one split at a time.

\end{document}